\documentclass[../main.tex]{subfiles}
\begin{document}

\conclusion{ДҮГНЭЛТ}

Энэхүү бакалаврын судалгааны ажлаар жижиглэн худалдааны салбарт тулгамдаж буй бараа материалын тооллого, аудитын асуудлыг шийдвэрлэх зорилгоор компьютерын хараа болон гүн сургалтын технологид суурилсан автомат систем хөгжүүлж, туршин үнэлсэн болно.

\section{Судалгааны үр дүнгийн нэгтгэл}

\subsection{Техникийн үр дүн}

Судалгааны хүрээнд дараах техникийн үр дүнд хүрсэн:

\begin{enumerate}
\item \textbf{Өндөр нарийвчлалтай объект илрүүлэлт:} YOLOv8-n загварыг 8 ангиллын бараа бүтээгдэхүүн дээр сургаж, mAP50 хэмжүүрээр 92.3\%, Precision 94.1\%, Recall 89.7\% нарийвчлалд хүрсэн. Энэ үзүүлэлт нь жижиглэн худалдааны орчинд бүрэн хэрэглэх боломжтой түвшин юм.
\item \textbf{Бодит хугацааны боловсруулалт:} Систем нэг зургийг дунджаар 245 миллисекундэд боловсруулж, секундэд 4 зураг боловсруулах чадвартай болсон. GPU ашиглан энэ хурдыг 28.5 FPS хүртэл нэмэгдүүлэх боломжтой.
\item \textbf{Автомат аудитын алгоритм:} Илрүүлсэн бараануудыг хүлээгдэж буй жагсаалттай харьцуулж, PASS/WARNING/FAIL статусыг 90\% нарийвчлалтай тодорхойлдог алгоритм боловсруулсан.
\item \textbf{Бүрэн бүтэн систем:} FastAPI backend, MongoDB database, мобайл болон вэб frontend-ээс бүрдсэн микросервис архитектурт систем хөгжүүлж, Docker-ээр контейнержуулсан.
\end{enumerate}

\subsection{Бизнесийн үр дүн}

Системийн бизнесийн хэрэглээний үр өгөөж дараах байдалтай байна:

\begin{enumerate}
\item \textbf{Цаг хугацааны хэмнэлт:} Гар аргаар 4--6 цаг зарцуулдаг тооллогыг 10--15 минутад багтаан хийх боломжтой болсон нь 80\%-ийн хэмнэлт юм.
\item \textbf{Алдааны бууралт:} Хүний хүчин зүйлээс үүдэлтэй 2--5\%-ийн алдааг 1\%-аас доош бууруулсан.
\item \textbf{Бодит хугацааны мэдээлэл:} Удирдлага дэлгүүрүүдийн нөхцөл байдлыг цаг алдалгүй хянах боломжтой болсон.
\item \textbf{Хямд өртөгтэй шийдэл:} Нэмэлт тоног төхөөрөмж шаардахгүй, зөвхөн ухаалаг утас ашиглан ажилладаг.
\end{enumerate}

\section{Судалгааны хувь нэмэр}

Энэхүү судалгааны ажлын шинжлэх ухаан, практикийн хувь нэмрүүд:

\begin{enumerate}
\item \textbf{Төгсгөл-хоорондын автомат систем:} Зураг авахаас эхлээд аудитын дүгнэлт хүртэлх бүх үйл явцыг автоматжуулсан цогц систем хөгжүүлсэн.
\item \textbf{Орон нутгийн датасет:} Монгол улсын зах зээл дээрх ундаа, усны бараануудын 324 зураг, 2,847 annotation-тай датасет бүрдүүлсэн.
\item \textbf{Практик хэрэглээ:} Академик судалгаанаас гадна бизнесийн бодит хэрэгцээнд шууд нэвтрүүлэх боломжтой систем бүтээсэн.
\item \textbf{Нээлттэй технологи:} Бүх бүрэлдэхүүн хэсгүүд нээлттэй эхийн технологи дээр суурилсан тул өргөтгөх, засварлах боломжтой.
\end{enumerate}

\section{Зорилтуудын биелэлт}

Судалгааны эхэнд дэвшүүлсэн зорилтуудын биелэлтийг үнэлэхэд:

\begin{table}[H]
\centering
\caption{Судалгааны зорилтуудын биелэлт}
\label{tab:objectives_completion}
\begin{tabular}{|p{5cm}|c|p{5.5cm}|}
\hline
\textbf{Зорилт} & \textbf{Биелэлт} & \textbf{Тайлбар} \\
\hline
Онолын судалгаа хийх & 100\% & YOLO хувилбаруудын харьцуулалт хийгдсэн \\
\hline
Өгөгдөл бүрдүүлэх ба загвар сургах & 100\% & 324 зураг, 8 ангилал, mAP50=92.3\% \\
\hline
Серверийн хөгжүүлэлт & 100\% & FastAPI + MongoDB бүрэн ажиллагаатай \\
\hline
Автомат аудитын алгоритм & 100\% & 90\% нарийвчлалтай ажиллаж байна \\
\hline
Гар утасны аппликейшн & 100\% & Odoo Mobile app iOS/Android дээр ажиллагаатай \\
\hline
\end{tabular}
\end{table}

\section{Хязгаарлалт ба сул тал}

Туршилтын явцад системийн дараах хязгаарлалтууд ажиглагдсан:

\begin{enumerate}
\item \textbf{Төстэй бараа андуурах:} Ижил төстэй шошготой боловч өөр амттай бараануудыг андуурах тохиолдол 2.3\% байсан.
\item \textbf{Гэрэлтүүлгийн нөлөө:} Бүдэг гэрэлтүүлэгтэй орчинд нарийвчлал 6\% орчим буурсан.
\item \textbf{Хагас халхлагдсан бараа:} 50\%-иас дээш халхлагдсан барааг илрүүлэх чадвар хязгаарлагдмал байв.
\item \textbf{Хязгаарлагдмал ангилал:} Загварыг зөвхөн 8 ангиллын бараан дээр сургасан.
\item \textbf{Интернет хамаарал:} Систем клоуд сервертэй холбогдож ажилладаг тул офлайн горимд ажиллах боломжгүй.
\end{enumerate}

\section{Цаашдын хөгжүүлэлтийн чиглэл}

Системийг цаашид дараах чиглэлээр өргөжүүлэн сайжруулах боломжтой:

\subsection{Богино хугацааны зорилтууд}

\begin{enumerate}
\item \textbf{Датасет өргөтгөл:} Өгөгдлийн санг 1000+ зураг, 50+ ангилал болгон өргөтгөх.
\item \textbf{Edge computing:} Гар утсан дээр шууд inference хийх боломжийг судлах.
\item \textbf{UI/UX сайжруулалт:} Хэрэглэгчийн санал хүсэлтэд үндэслэн интерфейсийг сайжруулах.
\end{enumerate}

\subsection{Дунд хугацааны зорилтууд}

\begin{enumerate}
\item \textbf{Video stream дэмжлэг:} Шууд камерын бичлэгээс таних боломжийг хөгжүүлэх.
\item \textbf{Planogram дэмжлэг:} Барааны байршлын зураглалтай уялдуулан аудит хийх.
\item \textbf{Multi-tenant архитектур:} SaaS платформ болгон хөгжүүлэх.
\end{enumerate}

\section{Эцсийн дүгнэлт}

Энэхүү бакалаврын судалгааны ажил нь орчин үеийн компьютерын харааны болон гүн сургалтын технологийг жижиглэн худалдааны салбарын бодит асуудлыг шийдвэрлэхэд хэрхэн амжилттай ашиглаж болохыг харуулсан практик жишээ болсон.

YOLOv8 загвар нь өндөр нарийвчлал, хурдан боловсруулалтын чадвараараа энэ төрлийн системд тохиромжтой сонголт болохыг баталсан. FastAPI болон Odoo-ийн хослол нь хөгжүүлэлтийг хурдасгаж, чанартай систем бүтээхэд тусалсан.

Хэрэгжүүлсэн систем нь жижиглэн худалдааны бизнес эрхлэгчид, түгээлтийн байгууллагуудад цаг хугацаа, хөдөлмөрөө ихээхэн хэмнэх боломжийг олгож байгаагаараа бодит ач холбогдолтой бүтээл болж чадсан.

\end{document}
